% !TEX program = xelatex
\documentclass{英语试卷}
\begin{document}
\examtitle{福建省 2025 届高中毕业班适应性练习卷}{英语}
\par 福建省 2025 届高中毕业班适应性练习卷\par
\par 高三年级 英语试题\par
\begin{notice}
\par 1. 答卷前，考生务必将自己的姓名、准考证号填写在答题卡上。\par
\par 2. 回答选择题时，选出每小题答案后，用铅笔把答题卡上对应题目的答案标号涂黑。如需改动，用橡皮擦干净后，再选涂其他答案标号。回答非选择题时，将答案写在答题卡上，写在本试卷上无效。\par
\par 3. 考试结束后，将本试卷和答题卡一并交回。\par
\end{notice}
\parttitle{第一部分 听力（共两节，满分 30 分）}
\par 做题时，先将答案标在试卷上，录音内容结束后，你将有两分钟的时间将试卷上的答案转涂到答题卡上。\par
\sectitle{第一节（共5小题：每小题1.5分，满分7.5分）}
\par 听下面 5 段对话。每段对话后有一个小题，从题中所给的 A、B、C 三个选项中选出最佳选项。听完每段对话后，你都有 10 秒钟的时间来回答有关小题和阅读下一小题。每段对话仅读一遍。\par
\question{1. Why do the speakers come here?}
\begin{choices}(3)
  \item To fetch breakfast.
  \item To have lunch.
  \item To cook supper.
\end{choices}
\question{2. What do we know about the delivery?}
\begin{choices}(3)
  \item It is free conditionally.
  \item It needs a whole day.
  \item It will be sent after 5 pm.
\end{choices}
\question{3. When will the flight probably arrive?}
\begin{choices}(3)
  \item At 6:00 pm.
  \item At 8:00 pm.
  \item At 9:00 pm.
\end{choices}
\question{4. What does the man want to do?}
\begin{choices}(3)
  \item Make a payment.
  \item Test a machine.
  \item Wash clothes.
\end{choices}
\question{5. What is the relationship between Peter and Sandra?}
\begin{choices}(3)
  \item Mother and son.
  \item Husband and wife.
  \item Father and daughter.
\end{choices}
\sectitle{第二节（共15小题：每小题1.5分，满分22.5分）}
\par 听下面 5 段对话或独白。每段对话或独白后有几个小题，从题中所给的 A、B、C 三个选项中选出最佳选项。听每段对话或独白前，你将有时间阅读各个小题，每小题 5 秒钟；听完后，各小题将给出 5 秒钟的作答时间。每段对话或独白读两遍。\par
\par 听第 6 段材料，回答第 6、7 题。\par
\question{6. Where did the woman go last month?}
\begin{choices}(3)
  \item A historical site.
  \item A silk museum.
  \item An antique shop.
\end{choices}
\question{7. What are the speakers mainly talking about?}
\begin{choices}(3)
  \item The history of the Silk Road.
  \item Ruins of the Great Wall.
  \item A travel to Yumenguan Pass.
\end{choices}
\par 听第 7 段材料，回答第 8 至 10 题，\par
\question{8. What suggestion does Mr. Smith offer Julie?}
\begin{choices}(3)
  \item Recording the lectures.
  \item Taking extra assignments.
  \item Working much harder.
\end{choices}
\question{9. Why does Julie consult Mr. Smith about the extra credits?}
\begin{choices}(3)
  \item She has missed some of his classes.
  \item She is worried about failing the course.
  \item She requires more tests for her grade.
\end{choices}
\question{10. When does the conversation probably take place?}
\begin{choices}(3)
  \item On Wednesday.
  \item On Thursday.
  \item On Friday.
\end{choices}
\par 听第8段材料，回答第11至13题。\par
\question{11. What does the man advise the woman to do?}
\begin{choices}(3)
  \item AI research.
  \item Data analysis.
  \item Voice acting.
\end{choices}
\question{12. What did the woman like as a little girl?}
\begin{choices}(3)
  \item Art.
  \item Biology.
  \item Computer.
\end{choices}
\question{13. What do we know about the woman's option of her particular future job?}
\begin{choices}(3)
  \item It is undecided.
  \item It is clear.
  \item It is disappointing.
\end{choices}
\par 听第9段材料，回答第14至17题。\par
\question{14. What do the speakers plan to do?}
\begin{choices}(3)
  \item Organize a game.
  \item Hold a BBQ party.
  \item Form a team.
\end{choices}
\question{15. Who are the basketball games intended for?}
\begin{choices}(3)
  \item Amateur players.
  \item Professional players.
  \item Girl players.
\end{choices}
\question{16. What problem do the speakers have to deal with?}
\begin{choices}(3)
  \item How to get enough girl participants.
  \item When to contact the organizer.
  \item Where to hold the event.
\end{choices}
\question{17. Who is a member of the school basketball team?}
\begin{choices}(3)
  \item Clark.
  \item Jane.
  \item Alice.
\end{choices}
\par 听第10段材料，回答第18至20题。\par
\question{18. Which was Jeffrey's best-seller?}
\begin{choices}(3)
  \item Gone with the Wind.
  \item The Art of Aging.
  \item The Last Lecture.
\end{choices}
\question{19. Where did Jeffrey work in 1980?}
\begin{choices}(3)
  \item In Detroit.
  \item In Orlando.
  \item In Chicago.
\end{choices}
\question{20. What did Jeffrey's column focus on?}
\begin{choices}(3)
  \item Famous people.
  \item Advice business.
  \item Social problems.
\end{choices}
\parttitle{第二部分 阅读理解（共两节，满分50分）}
\sectitle{第一节（共15小题：每小题2.5分，满分37.5分）}
\par 阅读下列短文，从每题所给的A、B、C和D四个选项中，选出最佳答案。\par
\passagelabel{A}
\par Here are four popular practice bases for secondary school students in New Zealand.\par
\par Southland Public Education Service\par
\par The Southland Public Education Service offers a variety of labor practice activities where students can learn new skills, make new friends and develop new interests, preparing them for a stable and successful future. The schools use flexible space designs to adapt to each student's unique learning style, enhancing creativity and efficiency. Students can participate in various interest groups and activities.\par
\par Greenfield Forest School\par
\par Greenfield Forest School provides cultural experiences, nature education, and outdoor adventures, where students can collaborate with environmental protection experts, wildlife scientists, and educators to explore real ecological environments and gain practical experience for upcoming challenges. The school offers seven thematic directions, including Wildlife Conservation Medical College, Climate Science Academy and so on.\par
\par Auckland Study Tour Program\par
\par The Auckland Study Tour Program aims to expand students' horizons and improve their English proficiency. Students can enhance their listening, speaking, reading, and writing skills through professional English teachers and explore the beautiful scenery of the region during their free time. The program helps students gain overseas living experience, laying a solid foundation for future study. The spring in New Zealand is perfect for nature lovers and outdoor enthusiasts.\par
\par New Zealand Study Program\par
\par The New Zealand Study Program integrates natural science, history, and technology, encouraging students to apply their knowledge comprehensively to solve practical problems and enhance interdisciplinary thinking. The program fosters students' cross-cultural communication skills and inclusiveness through multicultural environments and Maori cultural experiences, enhancing their global leadership abilities.\par
\question{21. What is special about Southland Public Education Service?}
\begin{choices}(2)
  \item It provides adaptable settings for individual needs.
  \item It offers practical cultural experiences for students.
  \item It introduces technology courses for integration.
  \item It organizes lively academic contests for nature lovers.
\end{choices}
\question{22. Which involves strengthening language competence?}
\begin{choices}(2)
  \item Southland Public Education Service.
  \item Greenfield Forest School.
  \item Auckland Study Tour Program.
  \item New Zealand Study Program.
\end{choices}
\question{23. What is shared by the four bases?}
\begin{choices}(2)
  \item Outdoor activities for scientific research.
  \item Cultural exchanges and global leadership.
  \item English training for daily communication.
  \item Skill development and future readiness.
\end{choices}
\passagelabel{B}
\par By the middle of the last decade, I had been studying Indian Butterflies for over 20 years. However, I had never observed butterflies at high altitudes in the trans-Himalayan habitat. So, in 2015, I decided to visit Ladakh to search for butterflies.\par
\par My first "butterflying" visit to Ladakh was in June 2016, then in July of the same year, followed by another visit in August 2018. For these visits, I charted out a two-week schedule on each occasion. I covered some of the well-known locations.\par
\par Each habitat we explored had its unique mix of butterfly species. By trickling streams and wet areas, we found Small and Dusky Green Underwings, along with numerous Meadowblues, all part of the Blues family. In alpine meadows, countless white and yellow butterflies from the Yellows and Whites family fluttered (翩翩飞舞) about. More than once, I'd spot a white or yellow butterfly flying across the meadow, and chase (追逐) after it, only to find myself out of breath in Ladakh's thin air while the butterfly gracefully flew away. Once, I chased a Butler's Dwarf butterfly for 15 minutes, managing only a few photos before it disappeared. Another time, I spotted dozens of Clouded Yellow butterflies in an alpine meadow. I spent half an hour chasing them, only for them to fly away each time.\par
\par Over the course of multiple visits to Ladakh for butterfly observation, I've had the remarkable opportunity to photograph more than 40 "lifers"—species I encountered for the very first time. These discoveries have added valuable knowledge to the field and led to the publication of two research papers in the "Journal of Threatened Taxa".\par
\par In addition to this, I partnered with the Snow Leopard Conservancy to create a small booklet showcasing 50 butterfly species found in Ladakh, with the goal of raising awareness and promoting butterfly watching as a unique and sustainable activity. This initiative aims to inspire both locals and tourists to explore the rich biodiversity of the region, while fostering a greater appreciation for Ladakh's delicate ecosystems.\par
\question{24. Why did the author decide to visit Ladakh?}
\begin{choices}(2)
  \item To survey butterflies at high altitudes.
  \item To search for butterfly habitats.
  \item To carry out a travel schedule.
  \item To recommend tourist attractions.
\end{choices}
\question{25. Which best describes the author's chasing butterflies?}
\begin{choices}(2)
  \item Challenging and fascinating.
  \item Scary and exciting.
  \item Embarrassing and amusing.
  \item Dangerous and interesting.
\end{choices}
\question{26. What is the author's purpose of making a booklet?}
\begin{choices}(2)
  \item To document insect species in India.
  \item To encourage eco-aware exploration.
  \item To promote local tourism in Ladakh.
  \item To compare butterfly research findings.
\end{choices}
\question{27. What is the best title for the text?}
\begin{choices}(2)
  \item Uniqueness of Butterflies
  \item A Fascinating Journey
  \item Biodiversity in Ladakh
  \item A Fluttering Discovery
\end{choices}
\passagelabel{C}
\par In today's fast-paced world, the art of mindful living has become increasingly essential. Mindfulness, the practice of being fully present in the moment, offers a pathway to inner peace and clarity among the chaos of modern life. It is not merely a technique but a way of being that transforms our relationship with ourselves and the world around us.\par
\par At its core, mindful living involves paying attention to our thoughts, feelings, and sensations without judgment. By observing our experiences with curiosity and openness, we can break free from automatic reactions and cultivate a deeper understanding of ourselves. This awareness allows us to respond to life's challenges with greater wisdom and kindness.\par
\par One of the key benefits of mindfulness is its ability to reduce pressure. When we focus on the present moment, we let go of worries about the past or future. Simple practices, such as mindful breathing or mindful walking, can root us in the here and now, creating a sense of calm and stability. Over time, these practices can rewire our brains, making us more resilient (有韧性的) to stress and challenges.\par
\par Moreover, mindful living enhances our relationships. By being fully present with others, we listen more deeply and communicate with greater effectiveness. This fosters empathy (共情) and connection, strengthening the bonds we share with family, friends, and colleagues. Mindfulness also encourages gratitude, helping us appreciate the small joys of life that often go unnoticed.\par
\par In essence, the art of mindful living is a journey of self-discovery and growth. It invites us to slow down, treasure each moment, and welcome life with a sense of wonder. As we cultivate mindfulness, we not only improve our own well-being but also contribute to a more kind and harmonious world. In a society that often values speed and productivity over presence and reflection, mindful living offers a timeless reminder of what truly matters.\par
\question{28. What can we know about mindfulness from the first two paragraphs?}
\begin{choices}(2)
  \item It is combination of life chaos with inner peace.
  \item It is awareness of the present without judgement.
  \item It is absence of observing our experiences.
  \item It is wisdom of understanding with openness.
\end{choices}
\question{29. What is the main idea of paragraph 3?}
\begin{choices}(2)
  \item Simple actions make a big difference.
  \item Calm and stability enhance life quality.
  \item Mindfulness eases stress and builds toughness.
  \item Mindfulness recalls the past and values the future.
\end{choices}
\question{30. Which is a practice of mindful living?}
\begin{choices}(2)
  \item Checking emails while socializing.
  \item Listening to music while doing homework.
  \item Enjoying each bite while having meals.
  \item Doing housework while watching a movie.
\end{choices}
\question{31. What is the author's overall attitude toward mindful living?}
\begin{choices}(2)
  \item Sympathetic.
  \item Reserved.
  \item Critical.
  \item Supportive.
\end{choices}
\passagelabel{D}
\par Neil Shubin's Ends of the Earth provides readers with a comprehensive overview of the geology, oceanography, glaciology (冰川学), geopolitics, and climatology of the planet's polar regions: Antarctica and the Arctic. Shubin writes clearly about various complex topics, integrating stories about his own fieldwork experiences in these places and arguing that polar science provides an angle to see the natural world and the extraordinary ways we have come to know it.\par
\par Shubin explains climate cycles' influence on global-scale glaciation. The "Snowball Earth" hypothesis (假说) suggests Earth was covered in ice twice: once 710 million years ago and a second time 660 million years ago. After each freezing period, Earth’s temperature rose quickly and the ice melted completely. These climatic changes affected the movement of early human ancestors, whose travels were influenced by habitat, glacial, and climate changes due to Earth’s orbital shifts.\par
\par Shubin notes that the story of the poles has largely been told by men who were explorers first and scientists second. To correct this, he highlights the work of three trailblazing women researchers in polar geology: Margaret Bradshaw, one of the first women to lead field teams in Antarctica; Maureen Raymo, who studies terrestrial (陆地上 的) weathering and the carbon cycle; and Mary Dawson, a paleontologist who discovered that the Arctic was warm 50 million years ago and animals could travel between continents.\par
\par While these regions seem remote from civilisation, Shubin reveals the impacts of pole changes on the other parts of the globe. These include increased greenhouse gas release, sped-up global warming, changes to animal populations, and the reemergence of diseases like anthrax. One of the most important impacts is on the Atlantic Meridional Overturning Circulation (AMOC), an ocean current that brings warm water to the north Atlantic and keeps Europe "warm" relative to other locations at the same latitude (纬度). A collapse of the AMOC, which some experts argue is already occurring, could bring colder weather to these countries.\par
\par "Having led my own explorations to polar regions for several decades, I feel the work has become a part of me," writes Shubin. "But the opposite is also true: our lives become part of the landscape itself." Following this line of reasoning, we find it important to remember that our lives in nonpolar regions affect what happens at the poles.\par
\question{32. What does Shubin's book aim to do?}
\begin{choices}(2)
  \item Share his own travel stories.
  \item Offer insights into polar science.
  \item Present tips for polar explorations.
  \item Discuss the geopolitics of the poles.
\end{choices}
\question{33. What does the underlined word "trailblazing" in paragraph 3 mean?}
\begin{choices}(2)
  \item Inspiring.
  \item Pioneering.
  \item Qualified.
  \item Devoted.
\end{choices}
\question{34. Why does the author mention "AMOC". in paragraph 4?}
\begin{choices}(2)
  \item To explain a theory.
  \item To clarify a concept.
  \item To make a prediction.
  \item To support an idea.
\end{choices}
\question{35. What is stressed in the last paragraph?}
\begin{choices}(2)
  \item Climate cycles restore polar balance.
  \item Polar changes escape human control.
  \item Human choices shape polar eco-system.
  \item Tourism growth outweighs polar protection.
\end{choices}
\sectitle{第二节（共5小题；每小题2.5分，满分12.5分）}
\par 阅读下面短文，从短文后的选项中选出能填入空白处的最佳选项。选项中有两项为多余选项。\par
\par Patterns are repeated behaviors that shape our daily lives. We follow them when we stick to the same morning routine, react similarly to stress, or face the same conflicts with others. \blanknum{36} The key is knowing which to keep and which to break.\par
\par Harmful patterns, like surfing through social media for hours or avoiding difficult tasks, form because our brains prefer easy and familiar actions. \blanknum{37} To break such cycles, start small: replace one bad habit with a positive activity, set clear daily goals, and ask friends to support your efforts.\par
\par Even good patterns can become problems if we never update them. A study method that worked in high school might not suit college courses. Staying too long with outdated routines limits our ability to adapt. \blanknum{38} People who refresh their habits regularly tend to solve problems faster and think more creatively. Try this: every few months, review your routines. Ask yourself, "Is this still working? Could there be a better way?" \blanknum{39}\par
\par As Walt Disney said, "Progress needs change." Breaking patterns feels uncomfortable at first, but staying stuck in unhelpful habits costs more over time. Begin with manageable steps: wake up 30 minutes earlier, try a new hobby, or write down your thoughts before making decisions. Each change weakens unhelpful patterns and builds confidence to tackle bigger challenges.\par
\par \blanknum{40} Keep those that help you grow, but don't let them become cages. By balancing useful routines with the courage to try new things, we create space for both stability and progress. Track your habits, stay curious, and note: small changes today can lead to big improvements tomorrow.\par
\begin{choices}(1)
  \item Patterns are very easy to follow.
  \item Remember: patterns are tools — not rules.
  \item While some patterns help us, others hold us back.
  \item In the end, breaking patterns creates space for balance.
  \item Sticking to old patterns risks missing better opportunities.
  \item Each time we repeat these behaviors, they become harder to change.
  \item Simple workspace changes and diverse dialogues stimulate new ideas.
\end{choices}
\parttitle{第三部分 语言知识运用（共两节，满分30分）}
\sectitle{第一节（共15小题；每小题1分，满分15分）}
\par 阅读下面短文，从短文后各题所给的A、B、C和D四个选项中，选出可以填入空白处的最佳选项。\par
\par Berger was walking by Lake Waban in Wellesley with his dog, Tommy. Just then, Tommy \blanknum{41} some birds landing on the icy lake. The ice \blanknum{42} when Tommy rushed toward the birds. Tommy fell into the freezing water.\par
\par Though \blanknum{43}, Berger tried to rescue his dog with a nearby boat. However, he also \blanknum{44} the icy lake. The Wellesley Fire Department \blanknum{45} arrived at the scene. Rescuers first pulled Berger out of the water, then turned their \blanknum{46} to Tommy.\par
\par Drone (无人机) footage recorded the dramatic \blanknum{} 47\blanknum{}. Following the most important rules for thin ice rescues, two firefighters approached the broken ice cautiously, lying \blanknum{} 48\blanknum{} on their stomachs to distribute their weight and \blanknum{} 49\blanknum{} further breaking. One rescuer extended a stick toward Tommy, while the other \blanknum{} 50\blanknum{} a safety rope. The dog, battling \blanknum{} 51\blanknum{} in the freezing water, managed to grab hold of the stick. Slowly and carefully, the firefighters pulled Tommy to \blanknum{} 52\blanknum{}.\par
\par The whole event showed Berger just how skilled and \blanknum{53} the firefighters were, and how important it is to stay calm and follow the right \blanknum{54} when there's an emergency. Berger learned a big lesson about how dangerous thin ice can be — and the rescue that saved them both showed him the strength of \blanknum{55} when times are tough.\par
\cz{41}{scared}{attracted}{spotted}{released}
\cz{42}{froze}{collapsed}{piled}{floated}
\cz{43}{thrilled}{surprised}{disappointed}{panicked}
\cz{44}{sank into}{rushed toward}{arrived at}{turned to}
\cz{45}{quietly}{quickly}{suddenly}{actively}
\cz{46}{attention}{trust}{comfort}{patience}
\cz{47}{change}{discovery}{operation}{influence}
\cz{48}{still}{calm}{tense}{flat}
\cz{49}{cause}{deepen}{avoid}{stimulate}
\cz{50}{cut}{secured}{shook}{removed}
\cz{51}{casually}{smoothly}{steadily}{desperately}
\cz{52}{safety}{peace}{shelter}{stability}
\cz{53}{hardworking}{sociable}{courageous}{generous}
\cz{54}{goals}{steps}{signs}{tracks}
\cz{55}{teamwork}{creativity}{optimism}{communication}
\sectitle{第二节（共10小题；每小题1.5分，满分15分）}
\par 阅读下面短文，在空白处填入 1 个适当的单词或括号内单词的正确形式。\par
\par A walk through the hutong near Beijing's Drum Tower is always like stepping back in time. But \blanknum{56} one recent cold winter night, during a walk through a well-traveled hutong there, I found myself suddenly \blanknum{57}(transport) back to my childhood years.\par
\par I was heading toward one of the \blanknum{58} (area) two excellent CD record stores, when a display in a shop window caught my eye. Upon \blanknum{59} (close) inspection, I realized the painting was a jigsaw puzzle, and the shop was devoted to such puzzles. When I stepped inside, the store's \blanknum{60} (toast) warmth relighted fires in my heart and memory. Long ago, while doing the puzzles and searching \blanknum{61} (endless) for the proper piece, I often found myself wondering whether perhaps I had lost a piece \blanknum{62} maybe my cat had stolen one. To effectively solve problems, we need effort and rest; when stuck, we need a break to bring clearness, finding the missing piece rapidly.\par
\par Among the valuable lessons and rewards that these puzzles offer \blanknum{63} (be) that patience and determination are enhanced. But what is most beneficial is how, like the detective Sherlock Holmes, we look for answers or clues, working tirelessly \blanknum{64} (find) the perfect fit. It's an approach \blanknum{65} pays off when we encounter any challenge that life presents.\par
\parttitle{第四部分 写作（共两节，满分 40 分）}
\sectitle{第一节（满分15分）}
\par 假定你是李华，你的英国笔友 David 来信询问中国学生如何在周末利用 AI 自主学习。请你用英语给他写一封回信，内容包括：\par
\par (1) 简要介绍;\par
\par (2) 分析利弊， 注意：\par
\par (1) 写作词数应为 80 左右;\par
\par (2) 请按如下格式在答题纸的相应位置作答。\par
\par\vspace{0.4em}
\begin{center}\begin{tabular}{{|l|p{\dimexpr\linewidth-3em-4\tabcolsep}|}}
\hline
Dear David, &  \\
 & Yours,Li Hua \\
\hline
\end{tabular}\end{center}
\vspace{0.4em}
\sectitle{第二节（满分25分）}
\par 阅读下面材料，根据其内容和所给段落开头语续写两段，使之构成一篇完整的短文。\par
\par Tom Perry concentrated hard. The rocket competition would begin shortly. Every detail mattered, and he was determined not to overlook a single one. The pressure was huge, but Tom was accustomed to being on top. He had never tasted defeat, and he wasn't about to start now.\par
\par The other children rarely approached him, and when they did, it felt forced. Tom preferred to keep to himself, confident in his ability to win yet again. He assembled (组装) his launch platform with great care, sorting through engines, igniters (点火器), and other tools in his toolbox. Suddenly, his igniter wire slipped from his grasp and disappeared into the grass. In a panic, he found another and carefully fitted it into place.\par
\par Tom checked his rocket again. It was perfect, but when he glanced at the others' rockets, he found lots of problems immediately. One rocket, however, had the look of a winner. And Tom's heart sank. A thin boy carried it carefully.\par
\par "Hi," a voice interrupted his focus. It was the thin boy, Ed, offering a nervous smile, "Nice rocket". Tom whispered thanks, eyes on his tools. He asked Ed where he tested his rockets because he had never seen him at the competitions before. Ed said he just practiced alone behind his house. Although Tom felt a bit nervous, he pretended not to care much about the contest. "Good luck", said Tom, and he continued to concentrate on his own rocket.\par
\par As Ed lifted his rocket, Tom felt an increasing anxiety. Just then, he noticed a fin (尾翼) on Ed's rocket had loosened. But Ed hadn't noticed it. Tom felt a bit sorry but decided not to say anything.\par
\par As Ed's turn approached, Tom's heart jumped. He had to speak up, despite the risk of losing. "Ed, hold up!" he shouted, drawing the attention of the crowd.\par
\par 注意:\par
\par (1) 续写词数应为 150 左右;\par
\par (2) 请按如下格式在答题卡的相应位置作答。\par
\writingpara{Paragraph 1: Ed stopped, turned around and looked at Tom confusedly. \blanknum{}}
\writinglines{8}
\writingpara{Paragraph 2: Tom felt the pain of defeat but also a sense of pride. \blanknum{}}
\writinglines{8}
\end{document}
